\section{Implementation and Comparison}
\label{sec:impl}
\begin{table}[!t]
\centering
\caption{Macro and controller summary}
\label{tab:impl}
\setlength{\tabcolsep}{4pt}
\begin{tabular}{@{}ll@{}}
\toprule
Technology & IHP SG13G2 130\,nm, BEOL HfAlO MEMRES\\
Array & 64$\times$64 1T1R, HV 300/450\,nm, pitch 2.26$\times$2.77\,\si{\micro\meter}\\
Macro & 226.5$\times$547.4\,\si{\micro\meter\squared}; 4096 MEMRES, 21\,930 MOS\\
Area split & array 21\,\%, col.\ drivers 21\,\%, WL 12\,\%, TDC 23\,\%\\

Controller & 21-state FSM, 12 ops, 1673 cells (0.025\,mm$^2$)\\
Signoff & 100\,MHz (slow corner); DRC, antenna and LVS clean\\
Energy & 65--105\,pJ per 512-bit access (0.13--0.2\,pJ/bit), slice SPICE\\
\bottomrule
\end{tabular}
\end{table}
\begin{table}[!t]
\centering
\caption{Comparison with RRAM and time-domain CIM macros}
\label{tab:prior}
\scriptsize
\setlength{\tabcolsep}{2pt}
\begin{tabular}{@{}llllll@{}}
\toprule
 & node & memory & readout & summed$^\mathrm{a}$ & reported metric\\
\midrule
\cite{chen2018rram65nm} & 65\,nm & 1\,Mb 1T1R & current SA, 3\,b & 9--25 & 14.8\,ns/cycle\\
\cite{wan2022neurram} & 130\,nm & 3\,Mb 1T1R & ADC, calibrated & 256 & 43\,TOPS/W (1\,b)\\
\cite{sayal2020timedomain} & 40\,nm & none on chip & counter TDC & 25--150, serial & 12.1\,TOPS/W\\
\cite{xue2024tdcim} & 180\,nm & 8\,kb SRAM & counter TDC & 32 per chain & 0.3--0.8\,TOPS/W$^\mathrm{b}$\\
this work & 130\,nm & 4\,kb 1T1R & 8-tap TDC, tunable & $\le8$ & 0.13--0.2\,pJ/bit$^\mathrm{c}$\\
\bottomrule
\multicolumn{6}{@{}p{\columnwidth}@{}}{\scriptsize $^\mathrm{a}$cells sensed in one access; $^\mathrm{b}$8$\times$8\,b at 1.2\,V; $^\mathrm{c}$slice SPICE, 0.45\,\si{\micro\second} per access; SA: sense amplifier.}\\
\end{tabular}

\end{table}
The macro is full custom, the controller synthesized and placed around it
(Table~\ref{tab:impl}). Against the MAC macros of Table~\ref{tab:prior}, with their
reference-tracked or calibrated readouts, this work trades throughput for a bitwise read whose
every timing and drive level is a register or a pad.
